% !TEX TS-program = lualatex+se
% !TEX encoding = UTF-8

\documentclass[14pt, a4paper]{extbook}

% --- USTAWIENIA STRONY ---
\usepackage{geometry}
\geometry{
	left=12mm,
	right=12mm,
	top=15mm,
	bottom=15mm,
	marginparwidth=0mm,
	marginparsep=0mm
	}
%\usepackage[letterpaper,showframe]{geometry}	
% usuwanie nagłówków/stopek z\,dokumentu
\usepackage[indent=0pt]{parskip} % ustawianie paragrafów
%\setstretch{1.5} % półtorej linii odległości między wierszami
\usepackage{titlesec} % ustawianie nagłówków
%\usepackage{fullpage} % redukcja marginesów
\usepackage{setspace} % odstępy między liniami
\usepackage{enumitem}% listy numerowane
% -----------------------------
\newcommand{\wstawplik}[1]{\input{/Users/Laser/Documents/Latex/include/#1}}%skraca kod
%----------- style nagłówków ----------------
\newcommand{\headingone}{\scshape\large}
\newcommand{\headingtwo}{\scshape\normalfont}
\newcommand{\margintext}[1]{\textit{\small #1}}
% przesunięcia poziome
\newcommand{\smallwidth}{\hskip 1em}
\newcommand{\medwidth}{\hskip 3pt}
\newcommand{\largewidth}{\hskip 12pt}
% przesunięcia pionowe
\newcommand{\linespace}{\vskip 7pt}
\newcommand{\paragraphspace}{\vskip 8pt}
\newcommand{\segmentspace}{\vskip 18pt}
\newcommand{\sectionspace}{\vskip 28pt}
%------------------------------------------

% --- GREGORIO ---	
\usepackage{gregoriotex}
\gresetlinecolor{gregoriocolor}
\gresetheadercapture{annotation}{greannotation}{string}% odczytuje adnotacje nad inicjałem jako ciąg, co pozwala na jego renderowanie
\gresetheadercapture{commentary}{grecommentary}{string}% odczytuje komentarz jako ciąg, co pozwala na jego renderowanie
\grechangedim{commentaryraise}{0.3cm}{scalable} % podnosi komentarz wystarczająco ponad górną linię
\grechangestyle{annotation}{\color{black}\small}% format adnotacji
\grechangestyle{commentary}{\color{gregoriocolor}\small}% format komentarza
\grechangestyle{abovelinestext}{\color{gregoriocolor}\small}% format napisu nad liniami
\grechangestyle{initial}{\fontsize{36}{36}\selectfont} % zmniejsza wielkość inicjału do bardziej proporcjonalnego rozmiaru (=36pt)
\grechangedim{intersyllablespacenotes}{0.18cm}{scalable}% Minimalna odległość między nutami różnych sylab.

% umieszczenie zawartości składu ich unikalnego argumentu po lewej i prawej stronie każdej linii
\let\grelocalleftbox\localleftbox
\let\grelocalrightbox\localrightbox

%\newcommand{\separator}{\greseparator{3}{40}}
%\newfontface\vectis{VectisMonumentalW01-Regular}
%\def\greinitialformat#1{\fontsize{20}{20}\vectis #1}
% ------------------

\usepackage[hyperindex=true]{hyperref}

%---- usun konflikt między paracol i luacolor ----
\wstawplik{conparlua}
%-----------------------------------------------

% --- SPIS TREŚCI -------
%\usepackage{imakeidx}
%\makeindex[title=Alfabetyczny spis program=texindy,options=-M styl.xdy]
%\usepackage{filecontents}
%\begin{filecontents*}{styl.xdy}
%;;; xindy style file
%(markup-locclass-list :open "\dotfill" :sep "")
%(define-attributes (( "textbf" "default" )) )
%(markup-locref   :attr  "textbf"     :open "\textbf{\hyperpage{" :close "}}")
%(markup-locref   :attr  "textit"     :open "\textit{\hyperpage{" :close "}}")
%(markup-locref   :attr  "textttt"     :open "\textttt{\hyperpage{" :close "}}")
%(markup-locref   :attr  "texttsc"     :open "\texttsc{\hyperpage{" :close "}}")
%(markup-locref   :attr  "default"     :open "\hyperpage{" :close "}")
%\end{filecontents*}
% -------------------------

% ---  Konfiguracja kolumn ---
\usepackage{paracol}
\setcolumnwidth{.49\textwidth, .49\textwidth}
\setlength{\columnseprule}{0.4pt}
\setlength{\columnsep}{.02\textwidth}
\colseprulecolor{red}% wymagany pakiet xcolor
\footnotelayout{m}
% ----------------------------

\makeatletter 
\patchcmd\pcol@hfil
 {\copy\pcol@tempboxa}
 {\luacolorresetboxcolor\pcol@tempboxa\copy\pcol@tempboxa}{}{\fail}
\makeatother
% ----------------------------

\footnotelayout{m}

%% --- UKŁAD DWUKOLUMNOWY Z TŁUMACZENIEM ---
%\usepackage{tcolorbox}
%\newtcolorbox{nutybox}{size=minimal, colback=white, colframe=red, arc=0mm, width=\linewidth, halign=left, valign=center}
%\newtcolorbox{tlumbox}{size=minimal, colback=white, colframe=red, arc=0mm, width=\linewidth, halign=left, valign=top, fontupper=\small\itshape} 
%\newcommand{\nuty}[1]{\begin{nutybox}#1 \end{nutybox}}
%\newcommand{\tlum}[1]{\begin{tlumbox}\tlumacz{#1} \end{tlumbox}}
%% ------------------------------------------------------

\usepackage{Tabbing}
\usepackage{longtable}
\usepackage[savepos]{zref}

% --- CZCIONKI DOKUMENTU ---
\usepackage[latin,polish]{babel}
\usepackage{fontspec}
%\usepackage{times}
%\usepackage[math]{anttor}
%\usepackage{tgtermes}
%\usepackage{mathptmx}
\usepackage{ebgaramond}
\usepackage{libertine}
\defaultfontfeatures{Ligatures=TeX}

%\setmainfont{TeX Gyre Termes}
%\setmainfont{TeX Gyre Heros}
%\setmainfont{Calluna Regular}
\setmainfont{EB Garamond}
%\setmainfont{Bodoni Moda}
%\newfontfamily{\sekfont}{Times New Roman}
\newfontfamily{\sekfont}{Libertinus Serif}
\newcommand\sekcjafont{\fontsize{16}{16}\selectfont \sekfont}

%----- Strona tytułowa -------------
\newfontfamily{\tytfont}{JSL Ancient}
\newcommand\tytulfont[1]{\fontsize{36}{30}\selectfont \tytfont#1\par}
\newcommand\podtytulfont[1]{\fontsize{18}{10}\selectfont \tytfont#1\par}

%\setmainfont[<font features>]{<name of font>} % sets the roman font
%\setsansfont[<font features>]{<name of font>} % sets the sans font
%\setmonofont[<font features>]{<name of font>} % sets the monospace font

% --- KONIEC CZCIONKI DOKUMENTU ---

% --- GRAFIKA I KOLORY ---
\usepackage{xcolor} % kolory
\usepackage{graphicx} % umożliwia import obrazów
\graphicspath{{img/}} % katalog obrazów
\usepackage{pgfornament}% ornamenty graficzne
% --- KONIEC GRAFIKA I KOLORY -----------

% --- KOMENDY WŁASNE ---
\newcommand{\red}[1]{\textcolor{red}{#1}}
\newcommand{\black}[1]{\textcolor{black}{#1}}
% --------------------------------
\newcommand{\tlumacz}[1]{{\small \textit{#1}}}
\newcommand{\inicjal}[2]{\textcolor{red}{\large\textbf{#1}}{#2}}
%--------------------------------
\titlespacing*{\chapter}{0pt}{0pt}{5pt}	
\titleformat{\chapter}[hang]{\filcenter\LARGE\scshape\bfseries}{}{0pt}{} % format nagłówków
\titlespacing*{\section}{0pt}{0pt}{0pt}	
\titleformat{\section}[hang]{\filcenter\headingone}{}{0pt}{} % format nagłówków
\titlespacing*{\subsection}{0pt}{0pt}{0pt}	
\titleformat{\subsection}[hang]{\filcenter\headingtwo}{}{0pt}{} % format podnagłówków

% ----------- RUBRYKA -------------------
\newcommand{\rubryka}[1]{
	\begin{spacing}{0.8}
	\begin{center}
		\parbox{0.9\textwidth}{\centering\textcolor{red}{\small\textit{#1}}}
	\end{center}
	\end{spacing}
	\linespace}
% ----------- SEKCJA --------------------
\newcommand{\sekcja}[2]{
	\begin{spacing}{0.6}
	\begin{center}
		\sectionspace
		\section*{#1}
		\begin{minipage}[c]{0.88\linewidth}
			\centering\red{\small\textit{#2}}
		\end{minipage}
	\end{center}	
	\end{spacing}
	\segmentspace}
% ----------- SUBSEKCJA --------------------	
\newcommand{\subsekcja}[1]{
	\begin{spacing}{0.8}
		\begin{center}
			\subsection*{#1}	
		\end{center}	
	\end{spacing}
	\paragraphspace}
%------------DEKOR ---------------------
\newcommand{\dekor}[2]{
	\begin{center}
		\vspace*{\fill}
		\includegraphics[width=#1\textwidth]{#2}
		\vspace*{\fill}
	\end{center}}
%------- Zarządzanie wersami w psalmach z tlumaczeniem -------------------------
\newcounter{licznik}
\newcommand{\licznikZero}{\setcounter{licznik}{0}}% zeruj licznik
\newcommand{\licznikUp}{\refstepcounter{licznik}}% zwiększ licznik 
\licznikZero
\newcommand{\BeginParallel}{\begin{spacing}{0.8} \begin{paracol}{2} \licznikZero}
\newcommand{\lat}[1]{\selectlanguage{latin}\ensurevspace{2\baselineskip}\begin{minipage}[t]{\linewidth}\ifnum \value{licznik}>0 \licznikUp \thelicznik .~#1 \else \licznikUp#1 \fi \end{minipage}\switchcolumn}
\newcommand{\latN}[1]{\selectlanguage{latin}\ensurevspace{2\baselineskip}\begin{minipage}[t]{\linewidth}#1\end{minipage}\switchcolumn}% kolumna bez numeracji wersów
\newcommand{\pol}[1]{\selectlanguage{polish}\begin{minipage}[t]{\linewidth}#1\end{minipage}\switchcolumn*}
\newcommand{\EndParallel}{\end{paracol} \end{spacing}}

% --- ZNAKI SPECJALNE ---
% 
%\newcommand †
% --------------------------

\begin{document}
\pagestyle{empty}% nagłówek i stopka pusta	
%\fontspec{times}
%\large
\setlength{\parindent}{0pt}
\gresetbracerendering{font}

\sekcja{Ad Vesperas}{Antiphonale Monasticum p.450}
\greannotation{Ant.}\greannotation{6. f} \gabcsnippet{(c4)AL(f)le(g')lu(h)ia,(f) (,)* al(gh)le(g')lu(f)ia,(d_c) (,) al(f)le(gh)lu(ge~)ia(f.) (::)}
\red{\large{Psalm 116} \linespace}
\greannotation{Psalm}\greannotation{116} \gabcsnippet{(c4) Lau(f)dá(gh)te(h) Dó(h)mi(h)num,(h) <b>om</b>(ixi)nes(h) <b>Gen</b>(g)tes:(h.) <c><v>\greheightstar</v></c>(:) lau(h)dá(h)te(h) e(h)um,(h) <i><ul>om</ul></i>(f)<i><ul>nes</ul></i>(gh) <b>pó</b>(g)pu(f)li:(f.) (::)
<c>2. </c>Quón(h)i(h)am(h) con(h)fir(h)má(h)ta(h) est(h) su(h)per(h) nos(h) mi(h)se(h)ri(h)<b>cór</b>(ixi)di(h)a(h) <b>e</b>(g)jus:(h.) <c><v>\greheightstar</v></c>(:) et(h) vé(h)ri(h)tas(h) Dó(h)mi(h)ni(h) ma(h)net(h) <i><ul>in</ul></i>(f) <i><ul>æ</ul></i>(gh)<b>tér</b>(g)num.(f.) (::)
<c>3. </c>Gló(h)ri(h)a(h) <b>Pa</b>(ixi)tri,(h) et(h) <b>Fí</b>(g)li(h)o,(h.) <c><v>\greheightstar</v></c>(:) et(h) Spi(h)rí(h)<i><ul>tu</ul></i>(f)<i><ul>i</ul></i>(gh) <b>Sanc</b>(g)to.(f.) (::)
<c>4. </c>Sic(h)ut(h) e(h)rat(h) in(h) prin(h)cí(h)pi(h)o,(h) et(h) <b>nunc</b>,(ixi) et(h) <b>sem</b>(g)per,(h.) <c><v>\greheightstar</v></c>(:) et(h) in(h) sǽ(h)cu(h)la(h) sæ(h)cu(h)<i><ul>ló</ul></i>(f)<i><ul>rum</ul></i>.(gh) <b>A</b>(g)men.(f.) (::)}
\linespace
\greannotation{Ant.}\greannotation{6. f} \gabcsnippet{(c4)AL(f)le(g')lu(h)ia,(f) (,) al(gh)le(g')lu(f)ia,(d_c) (,) al(f)le(gh)lu(ge~)ia(f.) (::)}
%% --- Magnificat ---ó á é ú í <c><v>\grecross</v></c> <c><v>\greheightstar</v></c
\sekcja{Magnificat}{}
\gresetinitiallines{1}\greannotation{Ant. Magn.}\greannotation{8. g}\gabcsnippet{(c3)VE(e!fwh)spe(f)re(f) au(fe~)tem(d) sáb(ef)ba(f)ti,(e.) <c><v>\greheightstar</v></c>(,) quæ(g) lu(h)cé(i_j)scit(i) in(i_j) pri(h')ma(f) sáb(h')ba(h)ti,(g.) (;)
ve(h)nit(h') Ma(h)rí(h)a(h) Mag(f')da(h)lé(g_f)ne,(e_') (,) et(e) ál(g')te(e)ra(f') Ma(g)rí(f.)a,(f.) (,) vi(f)dé(g_h)re(f_h) se(g)púl(fe~)crum,(d) al(ef)le(f)lú(e.)ia.(e.) (::)}
\linespace
\greannotation{Cant.}\gabcsnippet{(c3) Ma(e)gní(fe)fi(eh)cat(h.) <c><v>\greheightstar</v></c>(:) á(h)ni(h)ma(h) <i><ul>me</ul></i>(g)<i><ul>a</ul></i>(h) <b>Dó</b>(f)mi(e)num.(e.) (::)
<c>2. </c>Et(e) ex(f)sul(h)tá(h)vit(h) spí(h)ri(h)tus(h) <b>me</b>(i)us(h.) <c><v>\greheightstar</v></c>(:) in(h) De(h)o(h) sa(h)lu(h)<i><ul>tá</ul></i>(g)<i><ul>ri</ul></i>(h) <b>me</b>(f)o.(e.) (::)
<c>3. </c>Qui(e)a(f) re(h)spé(h)xit(h) hu(h)mi(h)li(h)tá(h)tem(h) an(h)cíl(h)læ(h) <b>su</b>(i)æ:(h.) <c><v>\greheightstar</v></c>(:) ec(h)ce(h) e(h)nim(h) ex(h) hoc(h) be(h)á(h)tam(h) me(h) di(h)cent(h) om(h)nes(h) ge(h)ne(h)<i><ul>ra</ul></i>(g)<i><ul>ti</ul></i>(h)<b>ó</b>(f)nes.(e.) (::)
<c>4. </c>Qui(e)a(f) fe(h)cit(h) mi(h)hi(h) ma(h)gna(h) qui(h) <b>pot</b>(i)ens(h) est:(h.) <c><v>\greheightstar</v></c>(:) et(h) sanc(h)tum(h) <i><ul>no</ul></i>(g)<i><ul>men</ul></i>(h) <b>e</b>(f)jus.(e.) (::)
<c>5. </c>Et(e) mi(f)se(h)ri(h)cór(h)di(h)a(h) e(h)jus(h) a(h) pro(h)gé(h)ni(h)e(h) in(h) pro(h)<b>gé</b>(i)ni(h)es(h.) <c><v>\greheightstar</v></c>(:) ti(h)mén(h)<i><ul>ti</ul></i>(g)<i><ul>bus</ul></i>(h) <b>e</b>(f)um.(e.) (::)
<c>6. </c>Fe(e)cit(f) pot(h)én(h)ti(h)am(h) in(h) brá(h)chi(h)o(h) <b>su</b>(i)o:(h.) <c><v>\greheightstar</v></c>(:) di(h)spér(h)sit(h) su(h)pér(h)bos(h) men(h)te(h) <i><ul>cor</ul></i>(g)<i><ul>dis</ul></i>(h) <b>su</b>(f)i.(e.) (::)
<c>7. </c>De(e)pó(f)su(h)it(h) pot(h)én(h)tes(h) de(h) <b>se</b>(i)de,(h.) <c><v>\greheightstar</v></c>(:) et(h) ex(h)al(h)<i><ul>tá</ul></i>(g)<i><ul>vit</ul></i>(h) <b>hú</b>(f)mi(e)les.(e.) (::)
<c>8. </c>E(e)su(f)ri(h)én(h)tes(h) im(h)plé(h)vit(h) <b>bo</b>(i)nis:(h.) <c><v>\greheightstar</v></c>(:) et(h) dí(h)vi(h)tes(h) di(h)mí(h)<i><ul>sit</ul></i>(g) <i><ul>in</ul></i>(h)<b>á</b>(f)nes.(e.) (::)
<c>9. </c>Su(e)scé(f)pit(h) Is(h)ra(h)ël(h) pú(h)e(h)rum(h) <b>su</b>(i)um,(h.) <c><v>\greheightstar</v></c>(:) re(h)cor(h)dá(h)tus(h) mi(h)se(h)ri(h)cór(h)<i><ul>di</ul></i>(g)<i><ul>æ</ul></i>(h) <b>su</b>(f)æ.(e.) (::)
<c>10. </c>Sic(e)ut(f) lo(h)cú(h)tus(h) est(h) ad(h) pa(h)tres(h) <b>nos</b>(i)tros,(h.) <c><v>\greheightstar</v></c>(:) A(h)bra(h)ham(h) et(h) sé(h)mi(h)ni(h) e(h)<i><ul>jus</ul></i>(g) <i><ul>in</ul></i>(h) <b>sǽ</b>(f)cu(e)la.(e.) (::)
<c>11. </c>Gló(e)ri(f)a(h) Pa(h)tri,(h) et(h) <b>Fí</b>(i)li(h)o,(h.) <c><v>\greheightstar</v></c>(:) et(h) Spi(h)rí(h)<i><ul>tu</ul></i>(g)<i><ul>i</ul></i>(h) <b>Sanc</b>(f)to.(e.) (::)
<c>12. </c>Sic(e)ut(f) e(h)rat(h) in(h) prin(h)cí(h)pi(h)o,(h) et(h) nunc,(h) et(h) <b>sem</b>(i)per,(h.) <c><v>\greheightstar</v></c>(:) et(h) in(h) sǽ(h)cu(h)la(h) sæ(h)cu(h)<i><ul>ló</ul></i>(g)<i><ul>rum</ul></i>.(h) <b>A</b>(f)men.(e.) (::)}
\linespace
\gresetinitiallines{1}\greannotation{Ant. Magn.}\greannotation{8. g}\gabcsnippet{(c3)VE(e!fwh)spe(f)re(f) au(fe~)tem(d) sáb(ef)ba(f)ti,(e.) (,) quæ(g) lu(h)cé(i_j)scit(i) in(i_j) pri(h')ma(f) sáb(h')ba(h)ti,(g.) (;)
ve(h)nit(h') Ma(h)rí(h)a(h) Mag(f')da(h)lé(g_f)ne,(e_') (,) et(e) ál(g')te(e)ra(f') Ma(g)rí(f.)a,(f.) (,) vi(f)dé(g_h)re(f_h) se(g)púl(fe~)crum,(d) al(ef)le(f)lú(e.)ia.(e.) (::)}
\sekcja{Oratio}{}
\wstawplik{dominus_vobiscum_solemnis}
\inicjal{O}{rémus.}\\
\inicjal{S}{píritum} nobis, Dómine, tuæ caritátis infúnde: \red{*} ut, quos sacraméntis paschálibus satiásti, tua fácias pietáte concórdes.\\
\inicjal{P}{er} Dóminum nostrum Jesum Christum, Fílium tuum: qui tecum vivit et regnat in unitáte ejúsdem Spíritus Sancti Deus, per ómnia sǽcula sæculórum.
\vspace{4mm}
\gresetinitiallines{0}\gabcsnippet{(c4) <c><sp>R/</sp>.</c> A(g.)men(gh..) (::)}
%% -----------------
\end{document}